\documentclass[10pt,a4paper]{amsart}
\usepackage[margin=19mm]{geometry}
\usepackage{amsmath,amssymb,mathtools}
\usepackage[scr=rsfs,cal=euler]{mathalfa}
\usepackage{xfrac}
\usepackage[unicode,hidelinks,bookmarks=false]{hyperref}
\newcommand{\ie}{i.e.\ }
\newcommand{\cf}{cf.\ }
\newcommand{\resp}{respectively\ }
\newcommand{\Subsection}{\subsection}
\providecommand{\nobreakdash}{\nobreak\hbox{-}}
\setlength{\emergencystretch}{2em}
\multlinegap0pt
 \newtheorem{theorem}{Theorem}
\newtheorem*{T1}{Theorem 1}
\newtheorem*{T2}{Theorem 2}
 \newtheorem{lemma}{Lemma}
 \newtheorem{corollary}{Corollary}
\newtheorem{proposition}{Proposition}
\newtheorem{definition}{Definition}
\newtheorem{example}{Example}
\newtheorem{assumption}{Assumption}
 \newtheorem{claim}{Claim}
 \newtheorem{remark}{Remark}
\newtheorem{observation}{Observation}
\newcommand{\A}{\mathbb A}
\newcommand{\E}{\mathbb E}
\newcommand{\mP}{\mathbb P}
\newcommand{\D}{\mathbb D}
\newcommand{\G}{\mathbb G}
\newcommand{\mH}{\mathbb H}
 \newcommand{\N}{\mathbb N}
 \newcommand{\T}{\mathbb T}
 \newcommand{\Z}{\mathbb Z}
 \newcommand{\R}{\mathbb R}
 \newcommand{\Q}{\mathbb Q}
 \newcommand{\mS}{\mathbb S}
\newcommand{\cA}{\mathcal A}
 \newcommand{\cC}{\mathcal C}
 \newcommand{\cD}{\mathcal D}
 \newcommand{\cK}{\mathcal K}
 \newcommand{\cW}{W^{1,1}}
 \newcommand{\cWD}{\cW_\cD}
 \newcommand{\cVD}{\cV_\cD}
 \newcommand{\cF}{\mathcal F}
\newcommand{\cG}{\mathcal G}
\newcommand{\cH}{\mathcal H}
  \newcommand{\cL}{\mathcal L}
\newcommand{\hcL}{\hat{\mathcal L}}
 \newcommand{\cM}{\mathcal M}
 \newcommand{\cN}{\mathcal N}
\newcommand{\cP}{\mathcal P}
\newcommand{\parts}{\mathscr{P}}
\newcommand{\cS}{\mathcal S}
\newcommand{\cT}{\mathcal T}
\newcommand{\cU}{\mathcal U}
\newcommand{\cV}{\mathcal V} 
\newcommand{\bq}{\mathbf{q}}
\newcommand{\hH}{\hat{H}}
\newcommand{\ovu}{\overline{u}}
\newcommand{\unu}{\underline{u}}
\newcommand{\ovK}{\overline{K}}
\newcommand{\ovX}{\overline{X}}
\newcommand{\unX}{\underline{X}}
\newcommand{\bh}{\overline{h}}
 \newcommand{\tH}{\tilde{H}}
\newcommand{\tL}{\tilde{L}}
\newcommand{\hC}{\hat{C}}
\newcommand{\bfh}{\mathbf{h}} 
\newcommand{\fM}{\mathfrak M}
\newcommand{\hu}{\hat{u}}
 \newcommand{\ovG}{\overline{G}}
 \newcommand{\bH}{\bar{H}}
 \newcommand{\bD}{\bar{D}}
\newcommand{\ub}{\bar{u}}
\newcommand{\bz}{\bar{z}}
\newcommand{\bw}{\bar{w}}
\newcommand{\bL}{\bar{L}}
\newcommand{\bM}{\bar{M}}
\newcommand{\bX}{\bar{X}}
 \newcommand{\af}{\alpha}
 \newcommand{\fui}{\varphi}
 \newcommand{\ep}{\varepsilon}
\newcommand{\be}{\beta}
\newcommand{\ga}{\gamma}
\newcommand{\bga}{\bar\gamma}
 \newcommand{\ka}{\kappa}
 \newcommand{\Ga}{\Gamma}
 \newcommand{\de}{\delta}
 \newcommand{\De}{\Delta}
 \newcommand{\om}{\omega}
 \newcommand{\Om}{\Omega}
  \newcommand{\la}{\lambda}
\newcommand{\La}{\Lambda}
 \newcommand{\te}{\theta}
 \newcommand{\si}{\sigma}
 \newcommand{\dga}{\dot\gamma}
 \newcommand{\lV}{\left\Vert}
 \newcommand{\rV}{\right\Vert}
 \newcommand{\rip}{\rangle}
 \newcommand{\lip}{\langle}
 \newcommand{\mL}{\mathbb L}
\renewcommand{\div}{\operatorname{div}}
\newcommand{\diver}{\operatorname{div}}
\newcommand{\rot}{\operatorname{rot}}
\newcommand{\flux}{\operatorname{flux}}
\newcommand{\grad}{\operatorname{grad}}
\newcommand{\lap}{\bigtriangleup}
\newcommand{\leb}{\operatorname{Leb}}
\newcommand{\Lip}{\operatorname{Lip}}
\newcommand{\dom}{\operatorname{dom}}
\newcommand{\diam}{\operatorname{diam}}
\newcommand{\sgn}{\operatorname{sin}}
\newcommand{\hes}{\operatorname{Hess}}
\newcommand{\Alt}{\operatorname{Alt}}
\newcommand{\entre}{\setminus}
\newcommand{\ann}{\operatorname{ann}}
\newcommand{\graf}{\operatorname{graph}}
\newcommand{\pr}{\operatorname{proj}}
\newcommand{\tr}{\operatorname{Tr}}
\newcommand{\sop}{\operatorname{supp}}
\newcommand{\dHp}{\dfrac{\partial H}{\partial p}}
\begin{document}
\title[Weak KAM theorems for subriemannian Lagrangians depending on]{Weak KAM theorems for subriemannian Lagrangians depending on the unknown function}
\author{Renato Iturriaga; Héctor Sánchez Morgado}
\date{}
\maketitle
\noindent\emph{Source and licence.} Weak KAM theorems for subriemannian Lagrangians depending on the unknown function. Version arXiv:2607.07966v3 (CC BY 4.0). This material is distributed under the Creative Commons Attribution 4.0 International licence, http://creativecommons.org/licenses/by/4.0/, as recorded in the arXiv OAI licence metadata for this exact version and shown on the abstract page https://arxiv.org/abs/2607.07966v3. Full rights notice: licenses/arxiv-2607.07966-CC-BY-4.0.txt.\par\smallskip
\noindent\textit{Editorial scope.} Counted unit: the original \texttt{theorem} environment \texttt{Cle} at source lines 309--315 together with its complete proof at lines 316--405; the delivered document keeps source lines 236--417 so that every referenced label and every citation resolves inside it. Native editable source is the paper's own concatenated TeX; the AMS wrapper is editorial formatting only. No publisher class, upstream script or unrelated artwork is redistributed.
 \section{Tonelli's Theorem for non-autonomous Lagrangians}
\label{sec:tonelli}
\begin{definition}\quad
  \begin{itemize} 
   \item  A family $\cF$  of curves $\ga:[a,b]\to M$ is absolutely
  equicontinuous if $\forall\ep>0$, $\exists\de>0$
  s.t. $\forall\ga\in\cF$ and any  disjoint subintervals $ ]a_1,b_1[,\cdots,]a_N,b_N[$ 
 of $[a,b]$:
  \[\sum_{i=1}^Nb_i-a_i<\de\implies \sum_{i=1}^Nd(\ga(b_i),\ga(a_i))<\ep\]
\item A
 family $\cG\subset L^1([a,b],m)$,  $m$ the Lebesgue measure on
 $[a,b]$, is uniformly integrable 
if given $\ep>0$ there is $\de>0$ such that 
\[f\in\cG, E\subset[a,b]\hbox{ measurable}, m(E)<\de\implies
  \int_E|f|<\ep.\]
   \end{itemize}
 
\end{definition}
\begin{remark}\label{basic}\quad
  \begin{enumerate}
  \item  A family $\cF\subset\cWD([a,b])$ such that $\{\|\dga\| : \ga\in\cF\}$
    is uniformly integrable, is absolutely equicontinuous.
  \item An absolutely equicontinuous family is equicontinuous.
  \item A uniform limit of absolutely equicontinuous  functions
    is absolutely continuous.
    \end{enumerate}
\end{remark}
\begin{theorem}{\cite{SM}}\label{hor}
  Let $(\ga_n)_n$ be a sequence in $\cWD([a,b])$ that converges uniformly to
  $\ga$, and such that
  $( \|\dga _n\|)_n$ is uniformly integrable.   Then $\ga\in\cWD([a,b])$.
  \end{theorem}

Let $L\in C(I\times\cD)$ be such that for each
$t\in I$, $x\in M$, the map $\cD_x\to\R$, $v\mapsto L(t,v)$ is twice
differentiable with $\partial^2L:I\times\cD\to\cD^*\otimes\cD^*$ 
continuous and the following properties hold
\begin{description}
\item [Uniform superlinearity] For all $k\ge 0$ there is $C(k)\in\R$.
  such that \[L(t,v)\ge k\|v\|+C(k)\hbox{ for all }(t,v)\in I\times\cD.\]
\item [Uniform boundedness] For all $r\ge 0$, we have
  \[A(r)=\sup\{L(t,v):\|v\|\le r\}<+\infty.\]
\item [Strict convexity] There is  $\ka>0$ such that for all  $t\in I$, $v,w\in\cD$ 
 \[\partial^2L (t,v)(w,w)\ge\ka\|w\|^2.\]  
\end{description}
For $\ga\in\cWD([a,b])$ we define the action
\[A_L(\ga)=\int_a^bL(t,\dga(t))\ dt.\]
\begin{lemma}\label{UI}
  For $c\in\R$ let
  $\cF(c)=\{\ga\in\cWD([a,b]): A_L(\ga)\le c\}.$
  Then the family
  $\{\|\dga(t)\|:\ga\in\cF(c)\}$ is uniformly integrable.
\end{lemma}
\begin{proof}
 By the uniform superlinearity,
      \[L(t,v)\ge K\|v\|+C(K).\]
      Given $\ep>0$ take $K>0$ such that 
  $$c-C(0)(b-a)<K\ep.$$

  Let $\ga\in\cF(c)$ and $E\subset[a,b]$ be measurable, then  
\[C(K)m(E)+K\int_E\|\dga(t)\|\ dt
   \le\int_E L(t,\dga(t))\ dt\]
and
   \[C(0)m([a,b]\entre E)\le \int_{[a,b]\entre E} L(t,\dga(t))\ dt \]
  Adding the inequalities we get 
  \[(C(K)-C(0))m(E)+C(0)(b-a)+K\int_E
    \|\dga(t)\|\ dt    \le A_L(\ga)\le c\]
which gives
\[\int_E\|\dga(t)\|\ dt\le\frac{c-C(0)(b-a)}K+
  \frac{C(0)-C(K)}Km(E) <\ep+\frac{C(0)-C(K)}{K}m(E).\]
This gives the uniform integrability of $\{\|\dga(t)\|:\ga\in\cF(c)\}$.
\end{proof}
  

\begin{theorem}\label{Cle}
   Let $(\ga_n)_n$ be a sequence in $\cWD([a,b])$ converging uniformly to
   $\ga$ such that
   \[\liminf_{n\to\infty}A_L(\ga_n)<+\infty.\]
   Then $\ga\in\cWD([a,b])$ and
\[A_L(\ga)\le \liminf_{n\to\infty}A_L(\ga_n).\]
\end{theorem}

\begin{proof}
  
Let $l=\liminf\limits_{n\to\infty}A_L(\ga_n)$,  extracting a subsequence
still denoted $\ga_n$, we have  $l=\lim\limits_{n\to\infty}A_L(\ga_n)$, and forgetting some of the first curves $\ga_n$, 
we can suppose that $\ga_n\in\cF(l+1)$ for all $n$. 
Lemma \ref{UI} implies that $(\|\dga_n\|)_n$ 
is uniformly integrable, and Theorem \ref{hor} that $\ga\in\cWD([a,b])$.
  
Let us show how we can reduce the proof to the case where 
$M$ is an open subset of $\R^d$,  $d =\dim M$ and the horizontal 
distribution is trivial. The lagrangian $L$ is bounded
below by $C(0)$ . If $[a',b']\subset[a,b]$, for all $n$ we have

\[A_L(\ga_n|[a',b'])\le
  A_L(\ga_n)-C(0)(b-b'+a'-a)\]
so that
\[\liminf_{n\to\infty}A_L(\ga_n|[a',b'])<+\infty.\]
Let now the partition $a_0 = 0 < a_1 < \ldots < a_p = 1$ and
$U_1, \ldots, U_p$ be such that
the horizontal distribution is trivial on $U_i$  and
$\ga([a_{i-1},a_i])\subset U_i$, $i=1,\ldots,p$,
It is enough to prove that  
\[A_L(\ga|[a_{i-1},a_i])\le \liminf_{n\to\infty}A_L(\ga_n|[a_{i-1},a_i])\]
because that implies and 
\begin{align*}
  A_L(\ga) =\sum_{i=1}^pA_L(\ga|[a_{i-1},a_i]) &\le
\sum_{i=1}^p\liminf _{n\to\infty} A_L(\ga_n|[a_{i-1},a_i])\\
&\le \liminf_{n\to\infty}\sum_{i=1}^pA_L(\ga_n|[a_{i-1},a_i])=\liminf_{n\to\infty}A_L(\ga_n)
\end{align*}
In the sequel we suppose that $M=U$ is an open set of $\R^d$, the
horizontal distribution is $U\times\R^m$ with $\langle ,\rangle$ the
euclidean inner product in $\R^m$, and that $\ga([a,b])$ and 
all $\ga_n([a,b])$ are contained in $U$.
We will write $\dga(t)=(\ga(t), \zeta(t))$, $\dga_n(t)=(\ga_n(t), \zeta_n(t))$.

   \begin{lemma}\label{LemCle}
     Let $U\subset\R^d$ be open, $L\in C([a,b]\times U\times\R^m)$ be
     twice differentiable in the the third variable with
     $\partial_{vv}L\in C([a,b]\times U\times\R^m,\R^{m^2})$, stricly convex
       and uniformly superlinear in the third variable and with
$\{L(t,x,v):\|v\|\le c\}$ bounded for each $c>0$.
     Given $K\subset U$ compact, $C>0$ and $\ep>0$, there exists $\de>0$
  such that if $x\in K$, $| x-y|\le\de$, $\|v\|\le C$ and $w\in\R^m$,
  then
  \[  L(t,y,w)\ge L(t,x,v)+\partial_v L(t,x,v)\cdot(w-v) -\ep.\]
  \end{lemma}
  For a fixed constant $C$ let
   \[E_C=\{t\in[a,b]:\|\zeta(t)\|\le C\}  \]
  Given $\ep> 0$, the compact set $K=\ga([a, b])\cup\bigcup_{n\in\N} \ga_n[a, b]$ 
and the  constant $C$ fixed above, we apply Lemma \ref{LemCle} to get $\de >0$
  satisfying its conclusion. By the compactness of $[a,b]$ and the
  continuity of $\ga$, 
  there is $\eta>0$ such that  $t\in[a,b]$, $d(x,\ga(t)) < \eta$
  imply $|x-\ga(t)| < \de.$
 Since  $\ga_n$ converges uniformly to $\ga$,
  there exists an integer $n_0$ such that, for each $n\ge n_0$
  we have $d(\ga_n(t),\ga(t)) < \eta$ for each $t\in[a, b]$.
Hence, for each $n\ge n_0$ and almost all $t \in E_C$, we have
  \[L (\ga_n(t),\zeta_n(t))\ge L (\ga(t), \zeta(t)) +
    \partial_v L(t,\ga(t), \zeta(t))\cdot(\zeta_n(t) -\zeta(t))-\ep,\]
 and from the uniform superlinearity, $L(t,\ga_n(t), \zeta_n(t))\ge C(0)$
 a. e.. Thus
  \begin{align}\nonumber
    A_L(\ga_n)&\ge \int_{E_C}L(t,\ga(t), \zeta(t))\ dt
                       +C(0)m([a,b]\entre E_C)\\
& +\int_{E_C}\partial_vL(t,\ga(t), \zeta(t))\cdot (\zeta_n(t) -\zeta(t))\ dt
\-\ep m(E_C) \label{1}
  \end{align}
  Since $\{\|\zeta_n(t)\|\}$ is uniformly integrable, $\zeta_n(t)$ converges to
  $\zeta(t)$ in the weak topology $\sigma(L^1,L^\infty)$.
Since $\|\zeta(t)\|\le C$ for $t\in E_C$, the function $\chi_{E_C}(t)\partial_vL(t,\ga(t), \zeta(t))$
is bounded. Thus
\[\int_{E_C}\partial_v L(t,\ga(t), \zeta(t))\cdot (\zeta_n(t) -\zeta(t))\to 0,\]
as $n\to\infty$. Taking limit in \eqref{1} we have
\[l=\lim_{n\to\infty}A_L(\ga_n)\ge \int_{E_C}L(t,\ga(t), \zeta(t))\ dt
  +C(0) m([a,b]\entre E_C)-\ep m(E_C).\]
Letting $\ep\to 0$ we have
\begin{equation}
  \label{eq1}
  l=\lim_{n\to\infty}A_L(\ga_n)\ge \int_{E_C}L(t,\ga(t), \zeta(t))\ dt
  +C(0) m([a,b]\entre E_C).
\end{equation}
Since $\zeta(t)$ is defined and finite for almost all $t\in[a,b]$ we
have that $E_C\nearrow E_\infty$ as $C\nearrow +\infty$ with
$m([a,b]\entre E_\infty)=0$. Since $L(t,\ga(t), \zeta(t))$ is bounded below by $C(0)$,
the monotone convergence theorem gives
\[\int_{E_C}L(t,\ga(t), \zeta(t))\ dt\to\int_a^bL(t,\ga(t), \zeta(t))\ dt \hbox{ as }C\to+\infty.\]
Letting $C\nearrow+\infty$ in \eqref{eq1} we finally obtain
\[l=\lim_{n\to\infty}A_L(\ga_n) \ge A_L(\ga)\]
\end{proof}

\begin{corollary}\label{low-semi}
  The action
  $A_L:\cWD([a,b])\to \R\cup\{+\infty\}$ is lower semicontinuous
  for the topology of uniform convergence in $\cWD([a,b])$.
\end{corollary}
\begin{proof}
  Let $\ga_n$ be a sequence in $\cWD([a,b])$ that converges
  uniformly to $\ga\in \cWD([a,b])$. We must show that
  \[\liminf_{n\to\infty}A_L(\ga_n)\ge A_L(\ga).\]
If $\liminf\limits_{n\to\infty}A_L(\ga_n)=+\infty$ there is nothing to prove.
In other case, the result follows from Theorem \ref{Cle}
\end{proof}

\begin{thebibliography}{99}
\bibitem[SM]{SM} \newblock H. S\'anchez Morgado. \newblock Homogenization for subriemannian Lagrangians. \newblock {\em Nonlinearity} {\bf 36} (2023) 3043--3067.
\end{thebibliography}
\end{document}
